\section{Conclusion}
Balanced 2-gap companions make global persistence deliberately uninformative: every parent always leaves $r-2$ children, so the complete-period population grows under protective, random, adversarial, and mixed selection alike. The local distinction is carried by conditional destruction relative to random, its cumulative hazard, and allocation. Finite real-sieve data instead supply observed cohort fractions, which require a separate transfer theorem.
The proof begins with the exact global recurrence and then replaces population counting by cumulative local hazard:
\begin{equation*}
\begin{aligned}
N_{k+1}&=(r_k-2)N_k,\\
w_r&=\frac{rf_r}{2},\\
D(Q)&=\sum_{r < Q}-\log(1-f_r),\\
P(Q)&=e^{-D(Q)}.
\end{aligned}
\end{equation*}
The main answer is that there is no maximum finite constant factor worse than random. If $w_r=w < \infty$, local survival decays only as $(\log Q)^{-2w}$; quadratic square-window supply still dominates, and the head probability series still diverges. Under the stated spatial and mixing premises, square windows are eventually nonempty and head 2-gaps recur infinitely often for every fixed finite $w$.
The first nontrivial boundary occurs when worsening grows with the filter. Under Section~7.1's cumulative quota/error conditions and the relevant spatial premises, the neutral exact-quota companion recovers the random survival scale. The biased exact-quota companion recovers the following general-hazard frontiers when its conditional effective skew follows the displayed schedule and its placement, availability, and mixing premises hold:
\begin{equation*}
\begin{aligned}
w_r=1+c\log r,\quad c < 1
&\Longrightarrow
\text{eventual square-window occupancy},\\
w_r=1+c\log r,\quad c < \frac12
&\Longrightarrow
\text{infinitely recurring head 2-gaps, with mixing}.
\end{aligned}
\end{equation*}
The allocation theorem explains why a percentage alone cannot locate a process in this phase diagram. Uniform, protective, and targeted allocation can apply the same adversarial budget and produce different local damage. The quantity that enters the theorem is therefore the conditional local hazard, not the policy label by itself.
The complete-period comparison is exact: every new filter destroys the fraction $2/r$ of expanded cyclic 2-gaps, so its cumulative excess over the neutral benchmark is zero. The two full-cycle diagrams expose the two parts of this statement separately. The destruction diagram makes the per-filter identity immediate; the survival diagram shows what compounding that identity does and how strongly it separates from the $c=1$ schedule. Neither diagram is weakened by being a direct rendering of the calculation: their value is that the local equality and cumulative consequence can each be checked visually in the representation best suited to it.
The finite square-window measurements then describe the remaining localization question. Through head $1129$, the observed population has mean ratio $0.967$ to the random expectation, and the measured one-step window rates through filter $19{,}429$ are at or below $2/r$. Fixed cohorts for $Q=17,101,251,503$ have signed effective coefficients between $-0.0353$ and $0.00908$; these are finite partial-cycle measurements around the exact cycle law. The measurements remain far from the $c=1$ window-failure scale, but they do not locate the distinguished head pair relative to the $c=1/2$ recurrence frontier.
The real sieve question is now expressible as a cumulative comparison rather than a vague claim of random or adversarial behavior: measure the observed cohort destruction $\widehat f_r$ generated by the CRT-coupled harmful indices, normalize it to $\widehat w_r$, and establish whether those observations control the relevant candidate hazards and placement near the head.
Here is one explicit sufficient deterministic counting criterion. Let $I_Q\in\{0,1\}$ indicate that the real sieve has a 2-gap at prime head $Q$. Let $\rho_Q>0$ be the companion reference weight obtained from the below-frontier cumulative hazard and the stated availability bound, and define
\begin{equation*}
R(X):=\sum_{\substack{Q\le X\\Q\text{ prime}}}\rho_Q.
\end{equation*}
Assume
\begin{equation*}
\begin{aligned}
R(X)&\longrightarrow\infty,
&&[\text{Divergent Reference Mass}]\\
\sum_{\substack{Q\le X\\Q\text{ prime}}}(I_Q-\rho_Q)
&=o(R(X)).
&&[\text{Deterministic Discrepancy Bound}]
\end{aligned}
\end{equation*}
Then
\begin{equation*}
\sum_{\substack{Q\le X\\Q\text{ prime}}}I_Q
=R(X)+o(R(X))
\longrightarrow\infty.
\qquad[\text{Q.E.D.}]
\end{equation*}
This is an unproved sufficient condition for the real CRT sieve. It is stronger than recurrence itself because it asserts an asymptotic counting formula. The reference mass must diverge, rather than every individual filter satisfying one pointwise bound. “Infinitely often” does not mean that every sufficiently large head is a 2-gap. Proving this condition would prove infinitely many head 2-gaps and therefore the twin-prime conjecture; it does not follow from the hazard estimates in this article.
\section{Future Work}
The companion theorems reduce the real-sieve question to measurable arithmetic inputs. The first direction is to estimate the real local hazard $D_{\kappa}(Q)$ from CRT-coupled filter locations and compare it with the head frontier without assigning hostile intent to individual filters. The second is to replace the stochastic mixing premise with a deterministic decorrelation or discrepancy theorem strong enough to transfer a divergent head-event sum to the real sequence. The third is to test exact-quota, delayed, block-balanced, and noisy-ranking companions on the same transitions, reporting both raw endpoint preference and quota-normalized effective skew.
These directions are deliberately separate. A finite experiment can identify which companion resembles the observed sieve, but it cannot establish the infinite transfer theorem. Conversely, a discrepancy theorem must control where the real CRT strikes land, not only how many strikes or global 2-gaps exist.
